\section{Cubic presentations of interacting coefficient operations}
\label{cp:sec}

The first-pole obstruction of Theorem~\ref{nl:thm:obstruction}
concerns a specified coefficient comparison over the normal base.
The formal order-atlas theory admits a different positive construction:
polynomial interactions can be resolved by cubic ones, and certain
changes of cubic presentation leave all transferred operations
unchanged.  The distinction between the coefficient systems remains
essential.

The simplest example is the quartic potential \(W=x^4\).  Over
\(k[\kappa]\), with \(|\kappa|=2\), its state differential is
\(d\eta=4\kappa x^3\).  The cubic potential
\[
 S_{\mathrm{small}}(x,u,p)=\kappa u^2+p(u-x^2),
 \qquad |x|=|u|=0,\quad |p|=2,
\]
replaces the quartic term by a quadratic constraint.  A canonical
construction uses three auxiliary pairs instead of one.  The two
contractions have the same original state complex as target, but
this fact alone compares no coefficient operation.  We will construct
a strict coefficient quotient and prove that it intertwines the
normalized homotopies.  It will identify every transferred tree map,
including its homotopy edges, on that original complex.

\subsection{Polynomial states and the coefficient domain}

Let \(R\) be a commutative, even-graded \(\mathbb Q\)-algebra with
zero differential.  Give finitely many even coordinates \(y_a\)
even integer degrees, and give their odd partners degree
\(|\eta^a|=1-|y_a|\).  A potential \(S(y)\) will have
cohomological degree two and ordinary polynomial degree at most three.
Its polynomial jet algebra is
\begin{equation}\label{cp:eq:states}
 E(S)=R[y_{a,r}:r\geq0]\otimes
                  \Lambda(\eta^a_r:r\geq0),
 \qquad
 dy_{a,r}=0,\qquad
 d\eta^a_r=\frac{T^r S_a}{r!},
\end{equation}
where \(S_a=\partial_aS(y_{1,0},\ldots,y_{N,0})\), and
\[
 Ty_{a,r}=(r+1)y_{a,r+1},\qquad
 T\eta^a_r=(r+1)\eta^a_{r+1},\qquad T|_R=0.
\]
The differential has degree one, squares to zero and commutes with
\(T\).  These assertions follow on generators because every \(S_a\)
contains only even closed coordinates.  Inputs are finite polynomials;
no completion of the input algebra is taken.

We use the coefficient algebra \(C_I\) of
Section~\ref{sec:free-state-trees}, extended to \(R\).
Explicitly, for a nonempty set \(J\) of total orders on a finite
label set \(I\), put
\[
 R_J=R[z_i:i\in I]
 \big[(z_i-z_j)^{-1}:
       \text{the comparison of \(i,j\) varies on \(J\)}\big].
\]
A coefficient is a face-compatible family in
\[
 R_J\otimes
 \mathbb Q[t_\sigma,dt_\sigma:\sigma\in J]/
       \bigl(\textstyle\sum t_\sigma-1,\sum dt_\sigma\bigr).
\]
The differential \(d_c\) sends \(t_\sigma\) to \(dt_\sigma\)
and kills \(R_J\).  Equivalently, one may use the full order simplex
and require each face restriction to lie in its stated overlap ring.
For distinct labels define
\begin{equation}\label{cp:eq:kernels}
 q_{ij}=\sum_\sigma t_\sigma\,1_{i<_\sigma j},
 \qquad
 c_{ij}=\frac{d_cq_{ij}}{z_i-z_j},
 \qquad
 c_{ij}^{r,s}=
 \frac{\partial_{z_i}^r\partial_{z_j}^s c_{ij}}{r!s!}.
\end{equation}
On a face with fixed comparison the fraction means zero without
inverting its denominator.  Otherwise that denominator is a unit.
Thus \(c_{ij}=c_{ji}\) is a closed degree-one coefficient, with
the same normalized jet kernels as
\eqref{fs:eq:order-kernel}.  Forgetting labels \(I\setminus J\)
gives the coefficient pullback \(C_J\to C_I\), not a specialization
of the omitted coordinates.

Write \(\mathcal T\) for a finite rooted insertion tree, reserving
\(T\) for translation.  At each vertex choose child positions and
an order simplex.  With internal scales \(\epsilon_e\), a leaf
coordinate is the sum of the successive displacements multiplied
by their ancestor scales.  Its order coordinates pull back by
\begin{equation}\label{cp:eq:flatten}
 t_\sigma\longmapsto
 \sum_{\operatorname{flat}(\sigma_v)=\sigma}
                     \prod_v t_{v,\sigma_v}.
\end{equation}
For fixed vertex faces, let \(B_{\mathcal T}\) be the tensor product
of their overlap rings and simplex-form algebras.  The scale ring is
\[
 B_{\mathcal T}[[\epsilon_e]]
                 [(\textstyle\prod_e\epsilon_e)^{-1}],
\]
as in \eqref{fs:eq:tree-ring}.  The admissible target
\(C_{\mathcal T}^{\mathrm{adm}}\widehat\otimes E(S)\) consists
of the coefficientwise expansions generated from a fixed polynomial
input, with one finite lower bound in every scale and finite
coefficients at each scale multiindex and auxiliary state weight.
For the latter weight give every state generator of jet index \(r\)
weight \(r+1\); this is separate from its cohomological degree.

The reason for the common pole bound is concrete.  If leaves \(i,j\)
first meet at a vertex \(v\), then
\begin{equation}\label{cp:eq:lca}
 Z_i-Z_j=s_v(u_a-u_b+r_{ij}),
\end{equation}
where \(s_v\) is the ancestor-scale monomial and \(r_{ij}\) has
positive descendant-scale degree.  Whenever the comparison varies,
\(u_a-u_b\) is inverted on the corresponding vertex overlap.
A pole \((z_i-z_j)^{-m_{ij}}\) therefore contributes only the
finite negative scale factor \(s_v^{-m_{ij}}\), followed by a
power series.  In particular, the lower bound at an edge \(e\) is
\[
 -\sum_{\{i,j\}:\ e\text{ lies above their first common vertex}}
                                                     m_{ij}.
\]
Refinement composes these substitutions on rational representatives:
the coordinate substitutions agree, and the inverse of each specified
unit is unique.  Expanding afterwards gives the same series.
There is no substitution on arbitrary unrelated negative Laurent
tails, nor an evaluation of every formal scale series at one.

\subsection{The cubic contacts and their strict insertion law}

The interaction needs an actual coefficient primitive.  For three
labels put \(\alpha=z_1-z_2\), \(\beta=z_2-z_3\) and
\(\gamma=z_1-z_3=\alpha+\beta\).  Enumerate the order vertices
as \(123,213,231,321,312,132\), and put
\(w_{ij}=t_i\,dt_j-t_j\,dt_i\).  Use
\begin{equation}\label{cp:eq:theta}
 \theta_{123}=
 -\frac{w_{02}+w_{35}}{2\alpha\gamma}
 +\frac{w_{04}-w_{13}}{2\beta\gamma}
 +\frac{w_{24}-w_{15}}{2\alpha\beta}.
\end{equation}
Each fraction is supported on a pair of orders which disagree on
its denominator pairs.  Its restriction to any other face is zero.
The formula is consequently a compatible form on the original atlas.
The transposition of labels \(1,2\) exchanges the first and third
summands with a minus sign and negates the second; the transposition
of \(2,3\) negates the first and exchanges the other two with a
minus sign.  Thus the formula is alternating, and equals its
normalized alternation.  Differentiating the Whitney forms gives
\begin{equation}\label{cp:eq:theta-d}
 d_c\theta_{123}
       =c_{12}c_{13}-c_{12}c_{23}+c_{13}c_{23}.
\end{equation}
Indeed \(d_cw_{ij}=2dt_i\,dt_j\); eliminating one \(dt_i\) with
\(\sum dt_i=0\) reduces the coefficient comparison to
\[
 \frac1{\alpha\gamma}-\frac1{\alpha\beta}
                       +\frac1{\beta\gamma}=0.
\]
Deletion defines the same primitive on any three labels of a larger
configuration.

One identity beyond \eqref{cp:eq:theta-d} controls insertion.
Insert \(1,2\) into one child and \(3\) into another.  If \(p\)
is the inner order coordinate and \(t\) the outer one, then
\[
 (t_0,t_1,t_2,t_3,t_4,t_5)
   =(pt,(1-p)t,0,(1-p)(1-t),p(1-t),0).
\]
Substitution in \eqref{cp:eq:theta} gives
\begin{equation}\label{cp:eq:theta-tree}
 \theta_{123}\longmapsto
 -\frac{(p-\tfrac12)\,dt}{(z_1-z_3)(z_2-z_3)}
 =(p-\tfrac12)\frac{c_{13}-c_{23}}{z_1-z_2}.
\end{equation}
For the first equality, \(w_{02}+w_{35}\) and \(w_{24}-w_{15}\)
pull back to zero, while \(w_{04}-w_{13}\) pulls back to
\((1-2p)\,dt\).  The second equality is the difference of the two
Cauchy kernels.  All normalized coordinate derivatives satisfy the
same identity.

Put \(Y_a(z)=e^{zT}y_{a,0}\).  For any translated even polynomial
\(H(z)=\sum_{n\geq0}H_nz^n\), its divided difference is the regular
series
\begin{equation}\label{cp:eq:divided}
 \frac{H(z)-H(w)}{z-w}
 =\sum_{n\geq1}H_n\sum_{\ell=0}^{n-1}z^{n-1-\ell}w^\ell.
\end{equation}
It introduces no inverse on an individual order chart.
The two cubic contact coefficients, before taking jet derivatives,
are
\begin{equation}\label{cp:eq:contacts}
 \begin{split}
 P^{(2)}_{ij;ab}
   &=(q_{ij}-\tfrac12)
       \frac{S_{ab}(Y(z_i))-S_{ab}(Y(z_j))}{z_i-z_j},\\
 P^{(3)}_{ijk;abc}&=S_{abc}\theta_{ijk}.
 \end{split}
\end{equation}
Here \(S_{ab}=\partial_a\partial_bS\) is affine linear, and
\(S_{abc}\) is a coefficient in \(R\).  Define
\(P^{(m)}_{i_1\cdots i_m;\mathbf a,\mathbf r}\) by applying
\(\prod_\nu\partial_{z_{i_\nu}}^{r_\nu}/r_\nu!\) to the corresponding
formula.  These derivatives act on the translated Hessian as well
as on the rational kernels.

All odd derivatives are left derivatives, and coefficient forms
are placed before state factors.  On distinct input copies set
\begin{equation}\label{cp:eq:wick}
 D_I=\sum_{i<j}\sum_{a,r,s}c_{ij}^{r,s}
 \left(
 \partial_{y_{a,r}^{(i)}}\partial_{\eta_s^{a(j)}}
 +\partial_{\eta_r^{a(i)}}\partial_{y_{a,s}^{(j)}}
 \right),
 \qquad
 \mathcal Y_I^0=\mathsf E_{\mathbf z}e^{D_I},
\end{equation}
where \(\mathsf E_{\mathbf z}\) translates each input by \(e^{z_iT}\)
and multiplies the results.  In characteristic zero this is precisely
the free product \(\mathsf E_{\mathbf z}\prod_{i<j}(1+D_{ij})\)
of \eqref{fs:eq:free-operation}: the pair operators commute and
each \(D_{ij}^2=0\).

For \(i_1<\cdots<i_m\) let
\[
 U_{i_1\cdots i_m;\mathbf a,\mathbf r}
 =\partial_{\eta_{r_m}^{a_m(i_m)}}\cdots
                  \partial_{\eta_{r_1}^{a_1(i_1)}}.
\]
This reversed derivative order sends the correspondingly ordered
product of odd generators to one.  On a homogeneous
coefficient-valued map \(F\), define
\begin{equation}\label{cp:eq:contact-operator}
 \mathscr A_m(F)=
 \sum_{\substack{i_1<\cdots<i_m\\\mathbf a,\mathbf r}}
 (-1)^{m|F|}
 P^{(m)}_{i_1\cdots i_m;\mathbf a,\mathbf r}
 F\,U_{i_1\cdots i_m;\mathbf a,\mathbf r},
 \qquad m=2,3.
\end{equation}
The coefficient multiplies the output of \(F\); the odd derivatives
act on its inputs.  They never differentiate a coefficient inserted
at the output by another application of \(\mathscr A_m\).
The full operation is
\begin{equation}\label{cp:eq:operation}
 \mathcal Y_I^S
      =\exp(\mathscr A_2+\mathscr A_3)
                         \mathsf E_{\mathbf z}\exp(D_I).
\end{equation}
Every exponential is finite on a polynomial input: a free contraction
removes one odd factor, and the two contacts remove two or three.
The remaining translations use the stated coefficientwise target.

\begin{proposition}[Strict cubic coefficient operations]
\label{cp:prop:cubic}
The maps \eqref{cp:eq:operation} are degree-zero chain maps on
\(E(S)\), and satisfy every finite rooted affine insertion identity
in \(C_{\mathcal T}^{\mathrm{adm}}\widehat\otimes E(S)\).
They preserve the unit, deletion, translations and relabellings.
If the even coordinates have spin \(s_a\), their odd partners spin
\(1-s_a\), and \(S\) has spin one, they are covariant for the
corresponding semilinear dilations.  A jet adds its index to spin.
\end{proposition}

\begin{proof}
The free pair has coefficient degree one and state-derivative degree
minus one.  In an \(m\)-contact the coefficient has degree
\(m-\sum_\nu|y_{a_\nu}|\), and the input derivatives have the
opposite degree.  Thus every term has total degree zero.

Let \(d_{\mathrm{Hom}}\) include both the coefficient differential
and the input and output state differentials.  Replacing each
\(P^{(m)}\) in \eqref{cp:eq:contact-operator} by \(d_cP^{(m)}\)
defines an operator \(\mathscr B_m\) of degree one.  The Gaussian
defect is
\[
 d_{\mathrm{Hom}}\mathcal Y_I^0
        =-(\mathscr B_2+\mathscr B_3)(\mathcal Y_I^0).
\]
To see every term, commute \(D_I\) with an even input coordinate.
The result is a constant odd-input derivative with an odd
coefficient, hence an even operator commuting with \(D_I\).
Conjugation therefore replaces that coordinate by itself plus this
operator in \(S_a\).  Taylor expansion of the quadratic polynomial
\(S_a\) gives only two- and three-input stars.  Their coefficients
are respectively
\[
 (S_{ab}(Y(z_i))-S_{ab}(Y(z_j)))c_{ij},
 \qquad
 S_{abc}(c_{ij}c_{ik}-c_{ij}c_{jk}+c_{ik}c_{jk}).
\]
Two edges from one centre to the same site have a repeated
\(d_cq\) and vanish.  Reversing the odd derivative order gives
the displayed star signs; the Taylor factorial cancels the orders
of its commuting contractions.  Remaining factors give
\(\mathcal Y_I^0U\).  Normalized coordinate differentiation gives
the same calculation for every jet, including derivatives of the
translated potential.

Equations \eqref{cp:eq:theta-d} and \eqref{cp:eq:divided} identify
these star coefficients with \(d_cP^{(m)}\).
The output coefficients involve only closed even fields, and \(U\)
supercommutes with the input differential.  The sign in
\eqref{cp:eq:contact-operator} consequently gives
\[
 [d_{\mathrm{Hom}},\mathscr A_2+\mathscr A_3]
       =\mathscr B_2+\mathscr B_3,
 \qquad
 [\mathscr A_2+\mathscr A_3,\mathscr B_2+\mathscr B_3]=0.
\]
The contact operators commute among themselves: exchanging their
coefficients and their odd derivative products contributes the same
sign twice.  Differentiating the finite exponential now cancels
the Gaussian defect.

For insertion, expand on a monomial input and distinguish repeated
factors temporarily.  A free contraction pairs one original even
factor with one original odd factor.  A binary contact consumes two
odd factors and creates at most one even output factor.  A ternary
contact consumes three odd factors and has scalar state output.
An outer free contraction can differentiate the even output of an
inner binary contact once.  It cannot differentiate that linear
output twice, and no contact creates an odd output.  Consequently
the connected contributions are exactly a free pair, an uncontracted
binary contact, and a ternary contact which may arise from one
binary contact followed by one free contraction.

Free pairs agree by the kernel Taylor identity
\eqref{fs:eq:kernel-translation}.  Binary contacts agree by regular
substitution in \eqref{cp:eq:divided}.  For a ternary contribution,
three leaves either lie in one child, lie in three different children
of their first common vertex, or split as two and one.  The first
case is decided in the child, and the second uses the three-flag
marginal of \eqref{cp:eq:flatten}.  In the last case the contracted
even output gives
\[
 S_{abc}(p-\tfrac12)
             \frac{c_{13}-c_{23}}{z_1-z_2},
\]
which is the direct ternary coefficient by
\eqref{cp:eq:theta-tree}.  These cases exhaust the connected
contributions in every finite tree.

Disjoint contributions exchange with two cancelling signs whenever
both the coefficient and the input operator are odd.  The exponential
factorials remove their possible orderings.  Restoring repeated
factors gives their polynomial derivative multiplicities in both
expansions.  Thus the entire coefficient-valued operations agree.
Normalized coordinate derivatives and the least-common-ancestor
bound \eqref{cp:eq:lca} justify every jet and pole expansion.
The order pullback is associative by flattening, so the identities
persist under every finite refinement.

A unit input is killed by every derivative, leaving the coefficient
deletion pullback.  The empty operation is \(1\), and the unary
operation is \(e^{zT}\), with identity at zero displacement.
Translation uses the exponential law.  The specified spin assignments
make the pole orders, contact coefficients and removed state factors
have the same dilation weights; normalized derivatives add their
jet indices on both sides.  Coefficients of nonzero spin transform
semilinearly.  Permutations preserve the kernels and the symmetric
potential derivatives, with the already specified tensor signs.
\end{proof}

\subsection{Polynomial cubic resolutions and normalized transfer}

Let \(k\) have characteristic zero, \(V\) be a finite-dimensional
species space, and \(W\) a polynomial on \(V\).  Retain the
integer-graded parameter ring \(k[\kappa]\), \(|\kappa|=2\), and
the original state complex
\begin{equation}\label{cp:eq:physical}
 A_W=k[\kappa,x_{a,r}:r\geq0]\otimes
                    \Lambda(\eta^a_r:r\geq0),
 \qquad
 d\eta^a_r=\frac{\kappa}{r!}T^r\partial_aW,\quad dx_{a,r}=0.
\end{equation}
Here \(|x|=0\), \(|\eta|=1\), and \(T\kappa=0\).
The parameter is neither inverted nor set to one.

Choose \(D\geq\max(2,\deg W)\), write
\(W(x)=\sum_{m=0}^D\langle W_m,x^m\rangle\), and normalize the
evaluation pairing by this equality.  Put \(u_0=1,u_1=x\).
For \(2\leq m\leq D\), introduce
\(u_m\in\Sym^mV\) of degree zero and
\(p_m\in(\Sym^mV)^*\) of degree two.  The potential
\begin{equation}\label{cp:eq:canonical}
 S_D=
 \kappa\sum_{m=0}^D\langle W_m,u_m\rangle
 +\sum_{m=2}^D\langle p_m,u_m-xu_{m-1}\rangle
\end{equation}
is cubic in the state coordinates.  Multiplication \(xu_{m-1}\)
is the symmetric product.  Define
\[
 v_m=u_m-xu_{m-1},\qquad
 u_m=x^m+\sum_{r=2}^m x^{m-r}v_r,
\]
and define \(H_r(x)\) by
\[
 \langle H_r(x),v\rangle
       =\kappa\sum_{m=r}^D\langle W_m,x^{m-r}v\rangle.
\]
With \(p'_r=p_r+H_r(x)\), direct substitution gives
\begin{equation}\label{cp:eq:split}
 S_D=\kappa W(x)+\sum_{r=2}^D\langle p'_r,v_r\rangle.
\end{equation}
The displayed triangular changes and their inverses are polynomial.
Every operation is symmetric multiplication or evaluation, so they
are natural under linear changes of species and coefficient extension.

Odd coordinates must be changed with them.  For an even coordinate
change \(y=y(z)\), use the full inverse Jacobian
\begin{equation}\label{cp:eq:odd-change}
 \eta'_b=\sum_a\frac{\partial y_a}{\partial z_b}\eta^a.
\end{equation}
Then \(d\eta'_b=\partial_{z_b}S\) by the chain rule.  Extend this
formula to jets by \(T^r/r!\).  The nonlinear changes used in
\eqref{cp:eq:split} are isomorphisms of state complexes.  The
coefficient operations on the split complex are defined by conjugating
the already specified operations on \(E(S_D)\); no nonlinear
normal-order covariance is assumed.

Each split auxiliary coordinate contributes a quartet
\[
 d\eta_{v,j}=p'_j,\qquad d\eta_{p',j}=v_j,\qquad
 dp'_j=dv_j=0.
\]
Let \(I:A_W\to E\) include the original split variables, and let
\(P:E\to A_W\) kill every auxiliary variable.  If \(N\) counts
the total polynomial and exterior degree in all auxiliary jets, put
\begin{equation}\label{cp:eq:homotopy}
 h=\frac1N\sum_{\mathrm{aux},j}
       \left(\eta_{v,j}\partial_{p'_j}
                    +\eta_{p',j}\partial_{v_j}\right)
 \quad(N>0),\qquad h=0\quad(N=0).
\end{equation}
The anticommutator of the auxiliary differential with the numerator
is its Euler operator \(N\); cross terms cancel by the odd signs.
The numerator squares to zero and preserves \(N\).  Hence
\begin{equation}\label{cp:eq:contraction}
 PI=1,\qquad dh+hd=1-IP,\qquad h^2=hI=Ph=0.
\end{equation}
Translation commutes with the numerator: in, for example,
\(\sum_j\eta_{v,j}\partial_{p'_j}\), the two translation
commutators cancel after shifting \(j\) by one.  It preserves \(N\),
so \(I,P,h\) commute with \(T\).  Only positive integers are
divided out.  Although \(I,P\) are unital state algebra maps,
they are not thereby strict coefficient maps.  Even the canonical
odd bracket detects this distinction: \(P\) kills an auxiliary
canonical pair whose bracket is one.

The normalized homotopy fixes all insertion coherences.
Let \(J\) be the dg algebra with basis \(e_0,e_1\) in degree zero
and \(a\) in degree minus one, and relations
\begin{equation}\label{cp:eq:interval}
 \begin{gathered}
 da=e_1-e_0,\qquad de_0=de_1=0,\\
 e_0^2=e_0,\qquad e_0e_1=e_1e_0=e_1^2=e_1,\\
 e_0a=ae_0=a,\qquad e_1a=ae_1=a^2=0.
 \end{gathered}
\end{equation}
The unit is \(e_0\).  The assignment
\begin{equation}\label{cp:eq:edge}
 \rho(e_0)=1_E,\qquad \rho(e_1)=IP,\qquad \rho(a)=-h
\end{equation}
is a dg algebra map to \(\operatorname{End}(E)\), by
\eqref{cp:eq:contraction}.  In particular
\(d_{\operatorname{End}}(-h)=IP-1_E\).

For a finite rooted tree \(\mathcal T\), put the strict cubic
operation at each vertex, \(I\) at each leaf, \(P\) at the root,
and \(\rho(\lambda_e)\) on each internal edge.  Use the coefficient
pullbacks \eqref{cp:eq:flatten} and the usual graded tensor rule.
This defines
\begin{equation}\label{cp:eq:tree-map}
 \mathcal M_{\mathcal T}:
 J^{\otimes E_{\mathrm{int}}(\mathcal T)}
     \otimes A_W^{\otimes L(\mathcal T)}
 \longrightarrow
 C_{\mathcal T}^{\mathrm{adm}}\widehat\otimes A_W.
\end{equation}
For a corolla it is \(P\mathcal Y_I^E I^{\otimes I}\).
Order the internal edges and put
\(\epsilon(e,\lambda)=\sum_{f<e}|\lambda_f|\).
Composition of the vertex, edge and leaf chain maps gives
\begin{equation}\label{cp:eq:cube}
 d_{\mathrm{Hom}}\mathcal M_{\mathcal T}(\lambda)
 =\sum_{\lambda_e=a}(-1)^{\epsilon(e,\lambda)}
 \left[
 \mathcal M_{\mathcal T}(\lambda[e:e_1])
 -\mathcal M_{\mathcal T}(\lambda[e:e_0])
 \right].
\end{equation}
An \(e_0\) face collapses its edge by strict cubic insertion.
An \(e_1\) face cuts that edge and composes the two transferred
trees, since its state operator is \(IP\).  Every face intersection
is an evaluation of this one tensor product of interval complexes;
its differential squares to zero with the displayed signs.
Deleting a unary identity vertex multiplies its adjacent decorations
in \(J\).  Grafting inserts an \(e_1\) edge.

These formulas are defined throughout the admissible target.
A polynomial leaf inclusion has finitely many factors.  Each vertex
removes odd factors, while each homotopy edge inserts at most one.
Thus a fixed tree and input permit finitely many contractions before
translation.  The coordinate changes have bounded polynomial degree,
and \(I,P,h\) commute with translation.  The normalized product rule
and \eqref{cp:eq:lca} therefore retain one finite pole bound
after every intermediate operation.  A homotopy divides no
coefficient parameter and introduces no new pole.

\subsection{A strict quotient for a quadratic constraint}

Return to a cubic \(S(y)\) over \(R\).  Let \(f_\alpha(y)\)
be finitely many degree-zero polynomials of ordinary degree at most
two.  Adjoin even \(u_\alpha,p_\alpha\) of degrees zero and two,
with odd partners \(\eta_u^\alpha,\eta_p^\alpha\), and put
\begin{equation}\label{cp:eq:extension}
 S^+(y,u,p)
   =S(y)+\sum_\alpha p_\alpha(u_\alpha-f_\alpha(y)).
\end{equation}
Let \(B\subset E(S^+)\) be generated by every jet of
\(y,\eta,p,\eta_u\), omitting every jet of \(u,\eta_p\).
On zero jets,
\[
 d\eta^a=S_a-\sum_\alpha p_\alpha\partial_a f_\alpha,
 \qquad d\eta_u^\alpha=p_\alpha.
\]
Thus \(B\) is a subcomplex.  Define \(q:B\to E(S)\) by killing
every jet of \(p,\eta_u\) and fixing the old coordinates.

\begin{theorem}[The strict constraint quotient]\label{cp:thm:quotient}
The inclusion \(B\hookrightarrow E(S^+)\) and the quotient
\(q:B\to E(S)\) are strict unital coefficient maps for the complete
operations \eqref{cp:eq:operation}.  They are quasi-isomorphisms
of their underlying state complexes, over \(R\) without localization.
\end{theorem}

\begin{proof}
Neither new free contraction acts on \(B\): the pairs
\(u_\alpha,\eta_u^\alpha\) and \(p_\alpha,\eta_p^\alpha\)
each require an omitted input generator.  Translation preserves the
omission.  In a nonzero contact, an input derivative can involve
only old \(\eta\)'s or \(\eta_u\)'s.  A potential derivative
involving \(p_\alpha\) requires an absent \(\eta_p^\alpha\)
input.  A derivative involving \(u_\alpha\) and no multiplier
vanishes beyond order one.  Hence the only additional Hessian
which contributes is
\[
 -\sum_\alpha p_\alpha\partial_a\partial_bf_\alpha
\]
on old coordinates.  Since each \(\partial_a\partial_bf_\alpha\)
is scalar, its binary divided difference and every normalized
coordinate derivative lie in the ideal of the \(p_\alpha\) jets.
The additional third derivatives on old coordinates vanish.

These facts show that each output stays in \(B\) and reduces
modulo \((p,\eta_u)\) to its old coefficient.  Crucially,
\eqref{cp:eq:contact-operator} differentiates inputs and multiplies
outputs.  Repeated contacts do not differentiate another inserted
output coefficient, so the same ideal argument applies term by
term to the entire finite exponential, including its free factors.
Consequently
\begin{equation}\label{cp:eq:strict}
 q\mathcal Y_I^{S^+}\big|_{B^{\otimes I}}
          =\mathcal Y_I^S q^{\otimes I}.
\end{equation}
The inclusion is strict because \(B\) carries the restricted
operation.  This proves strictness before any passage to cohomology.

For the state comparison set
\[
 v_\alpha=u_\alpha-f_\alpha(y),\qquad
 \xi_a=\eta^a+\sum_\alpha(\partial_a f_\alpha)\eta_u^\alpha.
\]
Then \(d\xi_a=S_a\), \(d\eta_u^\alpha=p_\alpha\), and
\(d\eta_p^\alpha=v_\alpha\).
The map \(j:E(S)\to B\) fixes \(y\) and sends \(\eta^a\)
to \(\xi_a\), with normalized jet extension; it is a chain
section of \(q\).  In these coordinates \(B\) is the old complex
tensored with the polynomial/exterior pairs \((p,\eta_u)\).
Their positive-count ideal is contracted by the corresponding
Euler homotopy.  The full \(E(S^+)\) has in addition the
contractible pairs \((v,\eta_p)\).  These explicit \(R\)-linear
splittings prove both quasi-isomorphisms, and survive arbitrary
coefficient extension.  No strictness of the section \(j\) is
asserted or needed.
\end{proof}

An invertible affine change of even coordinates, homogeneous for
the given degrees, is also strict after the inverse-transpose odd
change.  Its Jacobian and inverse are constant in the differentiated
state coordinates.  The free identity pairing is unchanged, while
the Hessian and third derivative transform by the cancelling
constant tensor powers.  Affine shifts cancel from Hessian
divided differences.  These statements hold after every normalized
jet derivative and prove strictness of \eqref{cp:eq:operation}.
They give no corresponding assertion for a nonlinear coordinate
change.

\begin{theorem}[Equality after normalized transfer]
\label{cp:thm:transfer}
Suppose \(E(S)\) is split as a state complex \(A\) with finitely
many auxiliary coordinate pairs and all their jets, and uses the normalized contraction
\eqref{cp:eq:homotopy}.  Split \eqref{cp:eq:extension} after that
old split, and normalize the contraction over all old and new
quartets together.  The transferred maps on \(A\) then satisfy
\[
 \mathcal M^+_{\mathcal T}(\lambda)
       =\mathcal M_{\mathcal T}(\lambda)
\]
for every finite rooted tree, every \(J\)-decoration, and every
polynomial input, as maps into the same admissible coefficient target.
\end{theorem}

\begin{proof}
In the split coordinates, \(B\) contains the old complex and the
new pairs \((p,\eta_u)\), while \(v,\eta_p\) are absent.
Let \(N_0\) be the old auxiliary count and \(N_1\) the new count.
The total homotopy preserves \(B\): its permitted new numerator
is \(\sum\eta_u\partial_p\), which never creates an omitted
variable.  Its old and new numerators preserve both counts.

If \(N_1>0\), \(q\) kills the input and the output of that
homotopy, so both \(qh^+\) and \(hq\) vanish.
If \(N_1=0<N_0\), the total denominator is \(N_0\), and
the homotopy is precisely the old one.  If both counts are zero,
both homotopies vanish.  Thus, in all cases,
\begin{equation}\label{cp:eq:intertwine}
 qh^+=hq,\qquad qI^+=I,\qquad P^+|_B=Pq,\qquad
 qI^+P^+|_B=IPq.
\end{equation}
For example \(h^+(p_{\alpha,j})=\eta_{u,j}^\alpha\) can be
nonzero when \(N_0=0\); it is the quotient of this output which
vanishes.  The count distinction is necessary.

Every leaf \(I^+\) lands in \(B\), every vertex preserves \(B\)
by Theorem~\ref{cp:thm:quotient}, and every edge operator
\(1,I^+P^+,-h^+\) preserves \(B\).  Starting with
\(P^+=Pq\) at the root, move \(q\) through each vertex using
\eqref{cp:eq:strict}, through each edge using
\eqref{cp:eq:intertwine}, and remove it at the leaves using
\(qI^+=I\).  The resulting expression is exactly the old
decorated-tree map.  The quotient has degree zero, so this
procedure changes no tensor sign.  It precedes all face evaluations,
and hence proves equality on every face intersection as well.
The quotient creates neither factors nor poles; the polynomial
sections and the old finite-tree bounds give the same admissible
completion throughout.
\end{proof}

\subsection{Canonical degree bounds and the quartic comparison}

Write \(\mathcal M^D\) for the transferred operations obtained
from \eqref{cp:eq:canonical} with its full split
\eqref{cp:eq:split} and normalized contraction.

\begin{corollary}[Independence of the canonical degree bound]
\label{cp:cor:degree}
For a fixed polynomial \(W\), integers
\(D'\geq D\geq\max(2,\deg W)\), and every decorated finite tree,
\[
 \mathcal M^{D'}_{\mathcal T}(\lambda)
                =\mathcal M^D_{\mathcal T}(\lambda)
\]
on \(A_W\).  The comparison and its inverse are identities.
All transition compositions for longer chains of degree bounds
are identities.  The equalities are natural under linear changes
of species and extension of the characteristic-zero coefficient
algebra, including specializations which lower the degree of \(W\).
\end{corollary}

\begin{proof}
Since \(W_{D+1}=0\), the next potential is exactly
\[
 S_{D+1}=S_D+
       \langle p_{D+1},u_{D+1}-xu_D\rangle.
\]
The symmetric-product constraint \(f=xu_D\) is quadratic and has
degree zero.  The old \(H_r\)'s are unchanged and \(H_{D+1}=0\).
Thus its graph split after the old split is the full canonical
triangular split, including its inverse-Jacobian odd change.
Its normalized total homotopy is the one used in
Theorem~\ref{cp:thm:transfer}.  That theorem gives equality for
one increase, and iteration gives every stated equality.
Symmetric multiplication, evaluation, normalized translation and
Euler counts commute with the specified coefficient and species
maps.  No exactness of an arbitrary tensor product is needed:
the identities and contractions themselves extend.
\end{proof}

The smaller quartic presentation requires two explicit steps.
Write \(u=u_2\), and compare
\begin{equation}\label{cp:eq:quartic-presentations}
 \begin{split}
 S_{\mathrm{can}}={}&\kappa u_4+
 p_2(u-x^2)+p_3(u_3-xu)+p_4(u_4-xu_3),\\
 S_{\mathrm{small}}={}&\kappa u^2+p(u-x^2).
 \end{split}
\end{equation}
First set \(q_4=p_4+\kappa\) and define
\[
 S_3=\kappa xu_3+p_2(u-x^2)+p_3(u_3-xu).
\]
Then make the invertible linear change
\(p=p_2-\kappa u,\ q_3=p_3+\kappa x\).
The exact potential identities are
\begin{equation}\label{cp:eq:quartic-steps}
 S_{\mathrm{can}}=S_3+q_4(u_4-xu_3),
 \qquad
 S_3=S_{\mathrm{small}}+q_3(u_3-xu).
\end{equation}
The inverse substitutions are
\[
 p_4=q_4-\kappa,\qquad
 p_2=p+\kappa u,\qquad p_3=q_3-\kappa x.
\]
Both are affine in the differentiated even coordinates over
\(k[\kappa]\), so they strictly identify the coefficient operations
by the affine argument above.  Each remaining step is a quadratic
constraint extension.

It remains to check the contractions; equality of the potentials
would not suffice.  For the small presentation put
\[
 v_2=u-x^2,\qquad
 p'_{\mathrm{small}}=p+2\kappa x^2+\kappa v_2.
\]
Then \(S_{\mathrm{small}}=\kappa x^4+p'_{\mathrm{small}}v_2\).
Its exact odd split is
\begin{equation}\label{cp:eq:small-odd}
 e_x=\eta_x+2x\eta_u-4\kappa x\eta_p,\qquad
 e_{v_2}=\eta_u-\kappa\eta_p,\qquad e_{p'}=\eta_p.
\end{equation}
The original differentials
\[
 d\eta_x=-2px,\qquad d\eta_u=2\kappa u+p,\qquad
 d\eta_p=u-x^2
\]
give \(de_x=4\kappa x^3,\ de_{v_2}=p'_{\mathrm{small}}\),
and \(de_{p'}=v_2\).  The small inclusion sends \(\eta\) to
\(e_x\).  Its projection sends
\[
 u\mapsto x^2,\qquad p\mapsto-2\kappa x^2,\qquad
 \eta_u,\eta_p\mapsto0,\qquad \eta_x\mapsto\eta.
\]
All jet components are their normalized translation extensions.

After the two graph splittings in \eqref{cp:eq:quartic-steps},
the full even coordinates are
\begin{equation}\label{cp:eq:quartic-full-split}
 \begin{gathered}
 v_2=u-x^2,\qquad v_3=u_3-xu,\qquad v_4=u_4-xu_3,\\
 p'_2=p_2+\kappa x^2,\qquad
 p'_3=p_3+\kappa x,\qquad p'_4=p_4+\kappa.
 \end{gathered}
\end{equation}
Indeed
\[
 p'_{\mathrm{small}}
 =p_2-\kappa u+2\kappa x^2+\kappa(u-x^2)
 =p_2+\kappa x^2.
\]
These are exactly the canonical coordinates of
\eqref{cp:eq:split}.  Their inverse is
\[
 \begin{gathered}
 u=x^2+v_2,\qquad
 u_3=x^3+xv_2+v_3,\qquad
 u_4=x^4+x^2v_2+xv_3+v_4,\\
 p_2=p'_2-\kappa x^2,\qquad
 p_3=p'_3-\kappa x,\qquad
 p_4=p'_4-\kappa.
 \end{gathered}
\]
Differentiating this entire inverse gives the physical odd coordinate
\begin{equation}\label{cp:eq:canonical-odd}
 \begin{split}
 e_x^{\mathrm{can}}={}&
 \eta_x+2x\eta_u+(u+2x^2)\eta_{u_3}
 +(u_3+xu+2x^3)\eta_{u_4}\\
 &-2\kappa x\eta_{p_2}-\kappa\eta_{p_3},
 \end{split}
\end{equation}
and the six auxiliary odd coordinates
\begin{equation}\label{cp:eq:canonical-aux-odd}
 \begin{aligned}
 e_{v_2}&=\eta_u+x\eta_{u_3}+x^2\eta_{u_4},
 &e_{p'_2}&=\eta_{p_2},\\
 e_{v_3}&=\eta_{u_3}+x\eta_{u_4},
 &e_{p'_3}&=\eta_{p_3},\\
 e_{v_4}&=\eta_{u_4},
 &e_{p'_4}&=\eta_{p_4}.
 \end{aligned}
\end{equation}
Their differentials are \(4\kappa x^3,p'_r,v_r\), respectively,
by the inverse-Jacobian chain rule.  The canonical inclusion sends
\(\eta\) to \eqref{cp:eq:canonical-odd}.  Its projection sends
\(u_r\) to \(x^r\), sends
\((p_2,p_3,p_4)\) to \((-\kappa x^2,-\kappa x,-\kappa)\),
kills the auxiliary odd coordinates, and sends \(\eta_x\) to \(\eta\).

The intermediate physical odd coordinate can also be followed through
the two steps.  The affine change from \(S_3\) to the second
expression in \eqref{cp:eq:quartic-steps} sends the odd coordinates
for \(x,u,p\) to
\[
 \eta_x-\kappa\eta_{p_3},\qquad
 \eta_u+\kappa\eta_{p_2},\qquad \eta_{p_2}.
\]
The graph section for \(u_3=xu\) adds
\(u\eta_{u_3}\) and \(x\eta_{u_3}\) to its first two entries.
Applying this section to \eqref{cp:eq:small-odd} gives
\begin{equation}\label{cp:eq:intermediate-odd}
 e_x^{(3)}=
 \eta_x+2x\eta_u+(u+2x^2)\eta_{u_3}
               -2\kappa x\eta_{p_2}-\kappa\eta_{p_3}.
\end{equation}
The next graph section, for \(u_4=xu_3\), adds
\((u_3+xu+2x^3)\eta_{u_4}\).  It gives exactly
\eqref{cp:eq:canonical-odd}, not a second choice of representatives.
The same calculation gives \eqref{cp:eq:canonical-aux-odd}.
It follows that both routes use the same full auxiliary counts
and the same normalized total Euler homotopy.

\begin{corollary}[The canonical and one-pair quartic theories]
\label{cp:cor:quartic}
With the coordinate maps and contractions just specified,
\[
 \mathcal M_{\mathcal T}^{\mathrm{can}}(\lambda)
       =\mathcal M_{\mathcal T}^{\mathrm{small}}(\lambda)
\]
on the original integer-graded \(A_{x^4}\), for every finite
decorated tree and every polynomial jet input.  The comparison and
inverse are the identity of \(A_{x^4}\).
\end{corollary}

\begin{proof}
Apply Theorem~\ref{cp:thm:transfer} first to the fourth constraint
and then to the third, using the strict affine identifications
\eqref{cp:eq:quartic-steps}.  Equations
\eqref{cp:eq:quartic-full-split}--\eqref{cp:eq:intermediate-odd}
identify the required normalized contractions with the original
ones in both presentations.  The two equalities of tree maps
therefore compose to the asserted equality.
\end{proof}

There are also strict structured arrows before transfer.
Let \(B_4\subset E(S_{\mathrm{can}})\) be the first constraint
subtheory, with strict quotient \(\pi_4:B_4\to E(S_3)\), and
let \(B_3\subset E(S_3)\) be the second, with quotient
\(\pi_3:B_3\to E(S_{\mathrm{small}})\).
The coefficient subtheory
\[
 C=\pi_4^{-1}(B_3)\subset B_4
\]
has strict maps
\[
 E(S_{\mathrm{can}})\longleftarrow C
                       \xrightarrow{\ \pi_3\pi_4\ }
                         E(S_{\mathrm{small}}).
\]
In split coordinates, \(C\) is \(B_3\) together with the ideal
generated by the fourth passive pair, with arbitrary \(E(S_3)\)
coefficients on that ideal.  Its positive fourth count is
contractible by \(\eta_{u_4}\partial_{q_4}/N_4\).
Together with the third passive-pair contraction and the omitted
quartet pairs, this proves that both displayed maps are
quasi-isomorphisms of state complexes.  Their strict coefficient
property follows from \(\pi_4\)'s strictness and the closure of
\(B_3\).  The normalized total homotopy preserves \(C\), because
it preserves each auxiliary group count and preserves \(B_3\).

\subsection{Presentation colours, the parameter and the remaining maps}

The equality of tree maps makes the coherent comparison explicit.
Put either presentation colour at each vertex, use the same
\(C_{\mathcal T}^{\mathrm{adm}}\) and the same insertion-edge
factors \(J^{\otimes E_{\mathrm{int}}(\mathcal T)}\), and assign
the common map \(\mathcal M_{\mathcal T}\), independently of
the colours.  A colour-changing unary vertex has state map
\(\id_{A_W}\).  Deleting such a vertex multiplies its adjacent
insertion-edge decorations in \(J\).

If an additional interval records a change of presentation, use a
separate copy \(J_{\mathrm{col}}\) with augmentation
\[
 \varepsilon_{\mathrm{col}}(e_0)
   =\varepsilon_{\mathrm{col}}(e_1)=1,\qquad
 \varepsilon_{\mathrm{col}}(a)=0.
\]
This respects the differential and multiplication of
\eqref{cp:eq:interval}.  Apply it to every added colour interval.
Both colour endpoints give the same map, and every positive-dimensional
colour cell gives zero.  The insertion-edge operator remains
\(\rho(a)=-h\); it is not replaced by zero.
Equation~\eqref{cp:eq:cube} still gives its complete signed
differential.  Edge collapse, grafting and all face intersections
therefore agree in both colours, as do all longer compositions of
the identity comparison and its identity inverse.

Units are retained throughout: the strict quotient fixes \(1\),
\(I(1)=1\), \(P(1)=1\) and \(h(1)=0\).
Derivatives at a unit input vanish, leaving exactly the existing
coefficient deletion map.  Relabelling permutes the same tensor
factors with the same Koszul signs.  Every coordinate map uses
normalized translation, and the Euler homotopies commute with it.
For the canonical and quartic presentations give \(x,u\) spin zero,
\(p,\kappa\) spin one, and the odd partners complementary spin.
All displayed affine and constraint formulas preserve these spins.
The dilation therefore retains its semilinear action
\(\kappa\mapsto s\kappa\); it is not asserted to be
\(k[\kappa]\)-linear at a fixed numerical coupling.

Parameter torsion is an immediate example of what this comparison
preserves.  At zero jets the quartic state complex has differential
\(d\eta=4\kappa x^3\).  Multiplication by \(4\kappa x^3\)
is injective on \(k[\kappa,x]\), so its cohomology is
\[
 k[\kappa,x]/(\kappa x^3)
\]
in even degrees and zero in odd degrees.
The class \(x^3\) is nonzero and killed by \(\kappa\).
Specializing the complex at \(\kappa=0\) restores the odd
generator.  Underived specialization of the cohomology ring alone
does not recover that generator.  The comparison and its contractions
use no inverse of \(\kappa\), so they preserve the original complex
and this distinction.

The equalities do not change the strict quartic defect of a different
primitive prescription.  To specify that defect, put
\[
 \Theta_{1234}=\operatorname{Alt}_{S_4}
 \bigl[(\theta_{123}-\theta_{124})c_{34}
             +(\theta_{134}-\theta_{234})c_{12}\bigr],
\]
where normalized alternation includes \(1/24\).  On the
disjoint-pair tree \(12\mid34\), let \(p,q\) be the two inner
order coordinates and \(r\) the outer one.  Set
\[
 F(\mathbf z)=
 \frac{z_1+z_2-z_3-z_4}
 {6(z_1-z_3)(z_1-z_4)(z_2-z_3)(z_2-z_4)}.
\]
The coefficient pullback is
\begin{equation}\label{cp:eq:old-quartic}
 \Theta_{1234}\longmapsto
 F(\mathbf z)\bigl((2q-1)\,dp+(2p-1)\,dq\bigr)\wedge dr.
\end{equation}
For example, grouping the alternating terms whose remaining edge
is \(12\) gives the \(dp\wedge dr\) coefficient
\[
 \frac{2q-1}{6(z_1-z_2)}
 \left[
 \frac1{(z_2-z_3)(z_2-z_4)}
 -\frac1{(z_1-z_3)(z_1-z_4)}
 \right]=(2q-1)F(\mathbf z).
\]
Exchanging the pairs gives the other coefficient.  A remaining
cross edge would contribute \(dr\wedge dr=0\), so these exhaust
the pullback.  It is nonzero, with allowed cross-pair poles, and
has primitive \(2F(\mathbf z)(p-\tfrac12)(q-\tfrac12)\,dr\).

For the quartic prescription with fourth contact
\(24\kappa\Theta_{1234}\), take four zero-jet odd inputs and
extract the coefficient linear in \(\kappa\), scalar in the state
variables, and of coefficient-form degree two.  The direct value
is \(24\kappa\Theta_{1234}\); the iterated binary value on
\(12\mid34\) is zero.  An inner binary contact would require two
outer free contractions to remove its quadratic even output, and
those contractions have the same outer \(dr\), whose square is
zero.  Two inner contacts have order \(\kappa^2\), while a root
binary contact leaves odd state factors.  There are no other terms
in that sector.  Thus the nonzero difference
\(24\kappa\Theta_{1234}\) and its allowed scalar primitive remain.
They are statements about this strict prescription, not the
transferred operations of Corollary~\ref{cp:cor:quartic}.

The comparison proved here covers canonical increases of \(D\),
the specified one-pair quartic presentation, and chains of the
verified quadratic extensions and compatible affine identifications
with their stated contractions.  It makes no assertion for an
arbitrary cubic presentation without such a chain.
Its operations are formal closed coefficient operations.
A supported action retaining the normal principal-part class of
Theorem~\ref{nl:thm:obstruction}, open-colour and module operations,
an associative chiral Koszul equivalence, analytic completion and
a physical state comparison still require their own maps.
The balanced-vacuum and Lie--chiral comparison of
Section~\ref{bv:sec} concerns a different state object and tensor.
Neither theorem identifies the coefficient construction here with
that vertex comparison.
